\subsection{训练监控与日志}

大规模训练中，监控指标是及早发现问题的唯一手段。一个完整的训练监控系统通常需要记录以下信息：

\begin{table}[H]
\centering
\begin{tabular}{p{0.20\textwidth}p{0.25\textwidth}p{0.35\textwidth}}
\toprule
\textbf{指标} & \textbf{正常行为} & \textbf{异常信号} \\
\midrule
Training loss & 单调下降（整体趋势） & 突然跳升、停滞不下、NaN \\
Learning rate & 按调度器变化 & 始终为 0 或异常大 \\
梯度范数 & 稳定在某个范围内 & 持续增大、突然飙升 \\
GPU 利用率 & 接近 100\% & 大幅波动或持续低于 50\% \\
显存占用 & 稳定不变 & 随步数缓慢增长（泄漏） \\
吞吐量 (tokens/s) & 稳定 & 周期性下降（可能是 checkpoint I/O） \\
\bottomrule
\end{tabular}
\caption{训练监控的核心指标与异常诊断}
\end{table}

\begin{lstlisting}[caption={训练循环中的日志记录示例}]
import wandb  # or tensorboard

wandb.init(project="cs336-lecture02")

for step in range(num_train_steps):
    x, y = get_batch(B=B)
    pred = model(x)
    loss = loss_fn(pred, y)
    loss.backward()

    grad_norm = torch.nn.utils.clip_grad_norm_(
        model.parameters(), max_norm=1.0
    )

    optimizer.step()
    optimizer.zero_grad(set_to_none=True)

    if step % log_interval == 0:
        wandb.log({
            "loss": loss.item(),
            "grad_norm": grad_norm.item(),
            "lr": optimizer.param_groups{[}0{]}{[}'lr'{]},
            "gpu_mem_gb": torch.cuda.memory_allocated() / 1e9,
        }, step=step)
\end{lstlisting}

\begin{knowledgebox}{Loss spike 的常见原因}
训练过程中偶尔出现的 loss spike（突然跳升）可能由以下原因引起：
\begin{itemize}
    \item \textbf{异常数据}：某个 batch 包含异常长的序列或罕见 token 组合
    \item \textbf{学习率过大}：尤其在 warmup 结束附近
    \item \textbf{数值不稳定}：梯度中出现 inf 或 nan
    \item \textbf{数据重复}：训练数据中存在大量重复导致过拟合后跳出
\end{itemize}
轻微的 loss spike 通常可以自行恢复，但如果 loss 持续不下降或出现 NaN，通常需要从上一个 checkpoint 回滚并降低学习率。
\end{knowledgebox}

